\documentclass[border=5pt]{standalone}
\usepackage{pgfplots}
\usepackage{pgf-pie} % 饼图：AI 代码可直接使用 \pie
\pgfplotsset{compat=1.18}
\usepgfplotslibrary{fillbetween}  % 面积图 / 堆叠面积图 / \closedcycle 填充路径
% 注：error bars 是 pgfplots 内置功能（在核心 pgfplots.errorbars.code.tex），不需要单独 \usepgfplotslibrary 加载
\usepackage{amsmath}
\usepackage{amssymb}
\usepackage{xcolor}

% 支持中文
\usepackage{fontspec}
\usepackage{xeCJK}
\setCJKmainfont{SimSun}
\setmainfont{Times New Roman}

\begin{document}

\begin{tikzpicture}
\begin{axis}[
    ybar,
    title={2024 年各省/直辖市 GDP 对比},
    ylabel={GDP（亿元）},
    symbolic x coords={广东,江苏,山东,浙江,四川,河南,湖北,福建,湖南,上海,安徽,河北,北京,陕西},
    xtick=data,
    xticklabel style={font=\scriptsize, rotate=45, anchor=east},
    nodes near coords,
    every node near coord/.append style={font=\scriptsize, yshift=3pt, fill=none, draw=none, inner sep=1pt, anchor=south},
    enlarge x limits=0.15,
    ymin=0,
    ymax=140000,
    legend style={at={(1.03,0.5)}, anchor=west, draw=black, fill=white},
    width=12cm,
    height=7cm
]
\addplot[fill=steelblue, draw=steelblue, thick] coordinates {
    (广东,135673)
    (江苏,128222)
    (山东,92069)
    (浙江,82553)
    (四川,60132)
    (河南,59132)
    (湖北,55803)
    (福建,54355)
    (湖南,50012)
    (上海,47218)
    (安徽,47050)
    (河北,43944)
    (北京,43760)
    (陕西,33786)
};
\legend{GDP}
\node[anchor=north west, font=\scriptsize] at (axis description cs:0.0,-0.15) {数据来源：2024 年各省统计公报};
\end{axis}
\end{tikzpicture}

\end{document}